% !PRINTED 661
\weberpage{661}
%<<CONT
eine Basis des K\"orpers $\Omega$, und folglich kann jede Funktion $\omega$
in $\Omega$ in der Form dargestellt werden:

\weq{3}{\omega = \xi_0 + \xi_1\theta + \cdots + \xi_{f-1}\theta^{f-1},}
worin die $\xi_0, \xi_1, \ldots, \xi_{f-1}$ rationale Funktionen von $z$ und
$z_1$ sind, also Zahlen des durch (1) bestimmten K\"orpers. Setzen wir

\weq{4}{n = ef, \quad n_1 = e_1 f,}
so lassen sich die Funktionen $\xi_i$ in jeder der beiden Formen
darstellen:

\weq{5}{\begin{aligned}
\xi &= x_0 + x_1 z_1 + x_2 z^2 + \cdots + x_{e-1} z_1^{e-1}, \\
% [sic] printed without subscript 1 on the third $z$; mathematically $z_1^2$ is intended
&= y_0 + y_1 z + y_2 z^2 + \cdots + y_{e_1-1} z^{e_1-1},
\end{aligned}}
worin die $x$ rationale Funktionen von $z$, die $y$ rationale Funktionen
von $z_1$ sind.

1. \gesperrt{Demnach l\"a\ss t sich jede Funktion $\omega$ linear mit
rationalen Koeffizienten in $z_1$ darstellen durch die $n_1$-Funktionen:}

\weq{6}{z^r\theta^s \quad \begin{aligned} r &= 0, 1, \ldots, e_1-1 \\ s &= 0, 1, \ldots, f-1 \end{aligned},}
\gesperrt{die wir in irgend einer Reihenfolge mit}

\weq{7}{\eta_1, \eta_2, \ldots, \eta_{n_1}}
\gesperrt{bezeichnen, und diese Funktionen sind linear unabh\"angig, d.\,h. es
besteht zwischen ihnen keine lineare Relation mit rationalen
Koeffizienten in $z_1$:}
\[
x_1\eta_1 + x_2\eta_2 + \cdots + x_{n_1}\eta_{n_1} = 0,
\]
\gesperrt{au\ss er wenn die $x_1, x_2, \ldots, x_{n_1}$ alle verschwinden.}

Denn $f$ ist der niedrigste Grad einer Gleichung f\"ur $\theta$ mit
rationalen Koeffizienten in $z$ und $z_1$ (\S\ 172).

Betrachten wir nun das Funktional

\weq{8}{\tau_1 = t_1\eta_1 + t_2\eta_2 + \cdots + t_{n_1}\eta_{n_1}}
mit den Funktionalvariablen $t_1, t_2, \ldots, t_{n_1}$, so ergibt sich
ein System von Gleichungen:

\weq{9}{\tau_1^r = x_{1,r}\eta_1 + x_{2,r}\eta_2 + \cdots + x_{n_1,r}\eta_{n_1},}
worin $r$ jede nat\"urliche Zahl, auch 0, sein kann, und die $x_{i,r}$
rationale holomorphe Funktionale in $z_1$ sind. Es ergibt sich daraus
durch Elimination der $\eta$ eine Gleichung

\weq{10}{\Phi(\tau_1) = \tau_1^m + A_1\tau_1^{m-1} + \cdots + A_{m-1}\tau_1 + A_m = 0}
h\"ochstens vom Grade $n_1$ mit rationalen Funktionalen in $z_1$ als
Koeffizienten.
